% ============================================================================
% Cover Letter — Paper 5 (JCAMD / Springer)
% Target: Journal of Computer-Aided Molecular Design (JCAMD)
% ============================================================================
\documentclass[11pt,a4paper]{letter}

\usepackage[margin=1.9cm]{geometry}
\setlength{\parskip}{0.3em}
\usepackage{microtype}
\usepackage{hyperref}
\usepackage{siunitx}
\sisetup{
  detect-all           = true,
  group-minimum-digits = 4,
}

\signature{Myke Vital Sao Temgoua \\ \small Corresponding Author}
\address{%
  University of Yaound\'{e} I \\
  Department of Physics, Faculty of Science \\
  P.O. Box 812, Yaound\'{e}, Cameroon \\
  \href{mailto:myke-vital.sao@facsciences-uy1.cm}{myke-vital.sao@facsciences-uy1.cm}
}

\begin{document}

\begin{letter}{%
  The Editors \\
  \textit{Journal of Computer-Aided Molecular Design} \\
  Springer Nature
}

\opening{Dear Editors,}

\noindent 29 August 2026

We submit the manuscript entitled \textbf{``Scaffold-Controlled Evaluation of Molecular Representations for Antimalarial Activity Prediction''} for consideration as a Research Article in the \textit{Journal of Computer-Aided Molecular Design}.

The work reports a controlled comparison of graph, sequence and topological molecular representations against an ECFP4--random-forest reference on a frozen panel of \num{19836} African antimalarial natural products. Evaluation proceeds under both stratified random and Bemis--Murcko scaffold cross-validation, each with five folds repeated over five seeds (\num{25} fold--seed replicates). The models compared are GIN, GIN--TFP (persistent-homology fusion), GIN--TNE (tensor-network fusion) and ChemBERTa.

Two protocol features are of particular relevance to JCAMD readers. First, pretrained ChemBERTa weights are restored independently before fine-tuning in every fold, eliminating the cross-fold weight-transfer channel that can otherwise inflate apparent transformer performance. This safeguard is rarely stated in molecular benchmarks yet is required for a valid transformer comparison. Second, the topological-fusion arms place persistent-homology and tensor-network descriptors inside a graph architecture and evaluate them under the same scaffold-held-out regime on an antimalarial natural-product panel---a combination not covered by recent representation-scale surveys. The primary ordering is further corroborated on a molecule-disjoint ChEMBL transfer panel and across three additional independent scaffold partitions; all artefacts are released in a versioned reproducibility package.

Under the scaffold protocol, ECFP4--RF attains the highest ROC AUC (\num{0.8300}), ahead of GIN--TFP (\num{0.8138}), GIN--TNE (\num{0.8090}), GIN (\num{0.8047}) and ChemBERTa (\num{0.7867}); all four learned arms remain significantly lower after Benjamini--Hochberg correction within the scaffold family. The same ECFP4--RF advantage over GIN is recovered on the external ChEMBL panel. Projection-weight analysis identifies persistent-image dimensions as the most salient block in the fusion head, yet this salience does not improve ranking. The result is presented as a data-set- and protocol-bounded negative finding, not as a universal claim about learned molecular models.

A phenotype-only LISH mechanism-of-action reference is supplied in the Supporting Information as orthogonal pharmacological context; its metrics are not numerically comparable with the molecular ROC-AUC results, and no structure-to-MoA claim is made. Code, frozen splits, trained checkpoints and the full Supporting Information are available in the public repository (\url{https://github.com/NanaEngo/Malaria_codesV2}); a citable Zenodo archive is being finalised and its DOI will be inserted at proof stage.

The work is original, is not under consideration elsewhere, has been approved by all authors, and the authors declare no competing financial interests.

\closing{We thank you for your consideration.}

\end{letter}
\end{document}
